\documentclass[10pt,a4paper]{amsart}
\usepackage[margin=19mm]{geometry}
\usepackage{amsmath,amssymb,mathtools}
\usepackage[scr=rsfs,cal=euler]{mathalfa}
\usepackage{xfrac}
\usepackage[unicode,hidelinks,bookmarks=false]{hyperref}
\newcommand{\ie}{i.e.\ }
\newcommand{\cf}{cf.\ }
\newcommand{\resp}{respectively\ }
\newcommand{\Subsection}{\subsection}
\providecommand{\nobreakdash}{\nobreak\hbox{-}}
\setlength{\emergencystretch}{2em}
\multlinegap0pt
\usepackage{bm}
\usepackage{bbm}
\usepackage{dsfont}
\usepackage{xspace}
\usepackage{cancel}
\usepackage{mathtools}
\usepackage{amsthm}
\makeatletter
\@ifundefined{lemma}{\newtheorem{lemma}{Lemma}}{}
\makeatother
\title[Complexity Framework For Forbidden Subgraphs V: Beyond Simpl]{Complexity Framework For Forbidden Subgraphs V: Beyond Simple Graphs}
\author{Tala Eagling-Vose, Barnaby Martin, Daniel Paulusma, Siani Smith}
\date{Source: arXiv:2502.07769v2 (CC BY 4.0)}
\begin{document}
\maketitle

\noindent\emph{Source and licence.} Complexity Framework For Forbidden Subgraphs V: Beyond Simple Graphs. Version arXiv:2502.07769v2 (CC BY 4.0), distributed under the Creative Commons Attribution 4.0 International licence, as recorded in the arXiv licence metadata for this exact version and shown on the abstract page https://arxiv.org/abs/2502.07769v2. Full rights notice: licenses/arxiv-2502.07769-CC-BY-4.0.txt.\par\smallskip

\noindent\textit{Editorial scope.} Counted unit: the Lemma whose source label is \texttt{lem:gen-obs} in the article (source0.tex lines 765--776, source label \texttt{lem:gen-obs}), delivered together with the article's own complete original proof of it (source0.tex lines 777--799, one \texttt{proof} environment of 23 lines). No mathematics is written, repaired or re-derived: statement and proof are the article's own text line for line. Required context retained uncounted: none. Every definition, display and label of those lines is retained unchanged. Omitted from this package and not redistributed: 1--764, 800--1290. Nothing inside a retained excerpt is omitted or altered: every retained body is delivered exactly as the article's lines are written, so no omission receipt of any kind accompanies this package. Citations: the retained excerpts carry no citation command at all, so no bibliography entry of the article is transcribed and no reference is redistributed. Reference adaptation: none was needed and none was made; every reference inside the delivered unit resolves inside it, so no unresolved reference or citation is produced. Printed slips kept literal: no printed word, symbol, hypothesis or number of the counted statement or of its proof was altered. Layout and class: the article's own class is \texttt{lipics-v2021} with packages including graphicx, tikz, amsmath, amssymb, color, xypic, todonotes, inputenc; the wrapper is the lane's pinned \texttt{amsart} editorial wrapper at 10pt on A4, which restates the theorem-like environments the retained text needs and adds the article's own macro definitions that the retained text needs (none), with extra wrapper packages (bm, bbm, dsfont, xspace, cancel, mathtools). The delivered numbering is the unit's own. Assets: no retained span contains a figure environment, a graphics-inclusion command, a PDF-page inclusion, a table or a TikZ picture; no image file is copied into this package, the package declares no asset dependency and no image file is redistributed with it. Licence: version arXiv:2502.07769v2 of this article is distributed under the Creative Commons Attribution 4.0 International licence, as recorded in the arXiv licence metadata for this exact version and shown on the abstract page https://arxiv.org/abs/2502.07769v2. A full rights notice is delivered at licenses/arxiv-2502.07769-CC-BY-4.0.txt. Characters: every retained excerpt is pure ASCII, so the delivered engine, XeLaTeX with its default Latin Modern text font, reports no missing character.
\begin{lemma}\label{lem:gen-obs}
There exists a polynomial time algorithm which, given an instance $G$ of {\sc Partially Reflexive Stable Cut}, either finds that $G$ is a yes-instance or returns an equivalent instance $G'$ with the following properties:
\begin{itemize}
    \item[1.] Every vertex of degree $2$ is reflexive and no vertex has degree ($0$ or) $1$.
    \item[2.] No degree $2$ vertex is contained in a triangle.
    \item[3.] There are no twin vertices.
    \item[4.] For any pair $u,v \in V$, if $N(u) \subseteq N(v)$, then $u$ must be reflexive and $v$ must not be.
\end{itemize}

Moreover, any simple graph contained as a subgraph in $G'$ is contained in $G$.

\end{lemma}

\begin{proof}
{\textbf {Property 1}}:
If a vertex $v$ of $G$ has degree ($0$ or) $1$ then either $N(v)$ is a stable cut of $G$ has a stable cut if and only if $G \setminus \{v\}$ does. Therefore, we either find that $G$ is a yes-instance or obtain an equivalent instance by deleting $v$.

Let $v$ be an irreflexive degree $2$ vertex. If $v$ is a cut vertex or $N(v)$ forms a stable cutset then $G$ is a yes-instance. Else, $G$ has a stable cutset if and only if $G \setminus {v}$ does and we may delete $v$ to obtain an equivalent instance. Repeating this process exhaustively, we either find that $G$ is a yes-instance or obtain an instance in which every vertex has degree at least $2$ and every degree $2$ vertex is reflexive.

\textbf{Property 2}:
Consider a degree $2$ vertex $v$ contained in a triangle. This vertex is reflexive by Property 1. Since the neighbourhood of $v$ is a clique, $N(v)$ is not a stable cut. Since $v$ is reflexive, it does not belong to any stable cut of $G$. Therefore, any stable cut of $G$ is a stable cut of $G \setminus v$. Since exactly one side of any stable cut in $G \setminus v$ contains a neighbour of $v$, $G$ has a stable cut if and only if $G \setminus v$ does and we may delete $v$. Repeating this process exhaustively, we obtain an instance in which every degree $2$ vertex has two independent neighbours.

\textbf{Property 3}:
Assume that $G$ contains a vertex $x$ with a true twin $y$. From Property $2$ we may assume $x$ has degree at least $3$. Assume that $G$ has a stable cut $C$ containing at least one of $x$ and $y$. It contains at most one of $x$ and $y$ since they are adjacent. Without loss of generality, $C$ contains $x$. Furthermore, since $x$ and $y$ have the same open neighboburhood, removing $x$ from $C$ does not connect the graph $G \setminus {C}$. Therefore, neither $x$ nor $y$ are contained in a minimal stable cut of $G$. Therefore, $G$ has a stable cut if and only if $G \setminus \{y \}$ has a stable cut which does not contain $x$. In other words, We obtain an equivalent instance by deleting $y$ and making $x$ reflexive.

We now assume that $x$ and $y$ are false twins. If there exists a stable cut of $G$ containing exactly one, say $x$, of $x$ and $y$ then, since $x$ and $y$ have the same open neighbourhood, there exists a stable cut containing neither of them. Therefore, any minimal stable cut $C$ either contains both $x$ and $y$ or neither of the two. In particular, if either of $x$ and $y$ is reflexive then we may make the other reflexive. Now, $G$ has a stable cut if and only if $G \setminus \{x \}$ does. In other words, we obtain an equivalent instance by deleting $x$.

Repeating these steps exhaustively, we obtain a graph with no twin vertices.

\textbf{Property 4}:
Note that $u$ and $v$ are non-adjacent since $v \notin N(v)$. Assume that $u$ is irreflexive. If $N(u)$ forms a stable cut then $G$ is a yes-instance, Since every neighbour of $u$ is a neighbour of $v$, $u$ is not a cut vertex. Therefore, any stable cut of $G$ is a stable cut of $G \setminus u $. On the other hand, consider a stable cut $C$ of $G \setminus u$. If $v$ belongs to $C$ then, since $u$ is irreflexive, we may add $u$ to $C$ to obtain a stable cut of $G$. If $v$ does not belong to $C$ then we add $u$ to the same side as $v$ to obtain a stable cut of $G$. Therefore, if $u$ is irreflexive then $G$ has a stable cut if and only if $G \setminus \{u\}$ does and we may delete $u$ to obtain an equivalent instance.

Now assume that both $u$ and $v$ are reflexive. Since every neighbour of $u$ is a neighbour of $v$, $G$ has a stable cut if and only if $G \setminus u$ does and we may delete $u$ to obtain an equivalent instance.

Finally, since each of these properties is obtained through vertex deletion, any simple subgraph contained in $G'$ is contained in $G$.%\qed
\end{proof}

\end{document}
